\documentclass[10pt,a4paper]{amsart}
\usepackage{amsmath,amssymb,amsthm,graphicx}
\usepackage{verbatim}
\newenvironment{reference}{\comment}{\endcomment}
%\newenvironment{reference}{}{}
\newenvironment{slogan}{\comment}{\endcomment}
\newenvironment{history}{\comment}{\endcomment}

% For commutative diagrams we use Xy-pic
\usepackage[all]{xy}

% We use 2cell for 2-commutative diagrams.
\xyoption{2cell}
% Theorem environments.
%
\theoremstyle{plain}
\newtheorem{theorem}[subsection]{Theorem}
\newtheorem{proposition}[subsection]{Proposition}
\newtheorem{lemma}[subsection]{Lemma}

\theoremstyle{definition}
\newtheorem{definition}[subsection]{Definition}
\newtheorem{example}[subsection]{Example}
\newtheorem{exercise}[subsection]{Exercise}
\newtheorem{situation}[subsection]{Situation}

\theoremstyle{remark}
\newtheorem{remark}[subsection]{Remark}
\newtheorem{remarks}[subsection]{Remarks}

\numberwithin{equation}{subsection}

% Macros
%
\def\lim{\mathop{\mathrm{lim}}\nolimits}
\def\colim{\mathop{\mathrm{colim}}\nolimits}
\def\Spec{\mathop{\mathrm{Spec}}}
\def\Hom{\mathop{\mathrm{Hom}}\nolimits}
\def\Ext{\mathop{\mathrm{Ext}}\nolimits}
\def\SheafHom{\mathop{\mathcal{H}\!\mathit{om}}\nolimits}
\def\SheafExt{\mathop{\mathcal{E}\!\mathit{xt}}\nolimits}
\def\Sch{\mathit{Sch}}
\def\Mor{\mathop{\mathrm{Mor}}\nolimits}
\def\Ob{\mathop{\mathrm{Ob}}\nolimits}
\def\Sh{\mathop{\mathit{Sh}}\nolimits}
\def\NL{\mathop{N\!L}\nolimits}
\def\CH{\mathop{\mathrm{CH}}\nolimits}
\def\proetale{{pro\text{-}\acute{e}tale}}
\def\etale{{\acute{e}tale}}
\def\QCoh{\mathit{QCoh}}
\def\Ker{\mathop{\mathrm{Ker}}}
\def\Im{\mathop{\mathrm{Im}}}
\def\Coker{\mathop{\mathrm{Coker}}}
\def\Coim{\mathop{\mathrm{Coim}}}

% Boxtimes
%
\DeclareMathSymbol{\boxtimes}{\mathbin}{AMSa}{"02}

%
% Macros for moduli stacks/spaces
%
\def\QCohstack{\mathcal{QC}\!\mathit{oh}}
\def\Cohstack{\mathcal{C}\!\mathit{oh}}
\def\Spacesstack{\mathcal{S}\!\mathit{paces}}
\def\Quotfunctor{\mathrm{Quot}}
\def\Hilbfunctor{\mathrm{Hilb}}
\def\Curvesstack{\mathcal{C}\!\mathit{urves}}
\def\Polarizedstack{\mathcal{P}\!\mathit{olarized}}
\def\Complexesstack{\mathcal{C}\!\mathit{omplexes}}
% \Pic is the operator that assigns to X its picard group, usage \Pic(X)
% \Picardstack_{X/B} denotes the Picard stack of X over B
% \Picardfunctor_{X/B} denotes the Picard functor of X over B
\def\Pic{\mathop{\mathrm{Pic}}\nolimits}
\def\Picardstack{\mathcal{P}\!\mathit{ic}}
\def\Picardfunctor{\mathrm{Pic}}
\def\Deformationcategory{\mathcal{D}\!\mathit{ef}}
\usepackage{cleveref}
\crefname{originallemma}{original lemma}{original lemmas}
\Crefname{originallemma}{Original lemma}{Original lemmas}
\makeatletter
\newlabel{derived-lemma-third-isomorphism-triangle}{{014A}{source}{}{sourcepointer-1}{}}
\newlabel{derived-lemma-third-isomorphism-triangle@cref}{{[originallemma][0][]014A}{[1][0][]source}{}{}{}}
\makeatother
\crefname{originalreference}{original reference}{original references}
\Crefname{originalreference}{Original reference}{Original references}
\makeatletter
\newlabel{lemma-derived-quasi-coherent-small-etale-site}{{071Q}{source}{}{sourcepointer-2}{}}
\newlabel{lemma-derived-quasi-coherent-small-etale-site@cref}{{[originalreference][0][]071Q}{[1][0][]source}{}{}{}}
\makeatother
\crefname{originalreference}{original reference}{original references}
\Crefname{originalreference}{Original reference}{Original references}
\makeatletter
\newlabel{lemma-induction-principle}{{08DR}{source}{}{sourcepointer-3}{}}
\newlabel{lemma-induction-principle@cref}{{[originalreference][0][]08DR}{[1][0][]source}{}{}{}}
\makeatother
\crefname{originalreference}{original reference}{original references}
\Crefname{originalreference}{Original reference}{Original references}
\makeatletter
\newlabel{lemma-single-complex-base-change-condition}{{0DJ8}{source}{}{sourcepointer-4}{}}
\newlabel{lemma-single-complex-base-change-condition@cref}{{[originalreference][0][]0DJ8}{[1][0][]source}{}{}{}}
\makeatother
\crefname{originalreference}{original reference}{original references}
\Crefname{originalreference}{Original reference}{Original references}
\makeatletter
\newlabel{lemma-unbounded-relative-mayer-vietoris}{{08HZ}{source}{}{sourcepointer-5}{}}
\newlabel{lemma-unbounded-relative-mayer-vietoris@cref}{{[originalreference][0][]08HZ}{[1][0][]source}{}{}{}}
\makeatother
\crefname{originallemma}{original lemma}{original lemmas}
\Crefname{originallemma}{Original lemma}{Original lemmas}
\makeatletter
\newlabel{perfect-lemma-single-complex-base-change}{{0DJ9}{source}{}{sourcepointer-6}{}}
\newlabel{perfect-lemma-single-complex-base-change@cref}{{[originallemma][0][]0DJ9}{[1][0][]source}{}{}{}}
\makeatother
\crefname{originallemma}{original lemma}{original lemmas}
\Crefname{originallemma}{Original lemma}{Original lemmas}
\makeatletter
\newlabel{perfect-lemma-single-complex-base-change-condition}{{0DJ8}{source}{}{sourcepointer-7}{}}
\newlabel{perfect-lemma-single-complex-base-change-condition@cref}{{[originallemma][0][]0DJ8}{[1][0][]source}{}{}{}}
\makeatother
\crefname{originalreference}{original reference}{original references}
\Crefname{originalreference}{Original reference}{Original references}
\makeatletter
\newlabel{remark-match-total-direct-images}{{08GH}{source}{}{sourcepointer-8}{}}
\newlabel{remark-match-total-direct-images@cref}{{[originalreference][0][]08GH}{[1][0][]source}{}{}{}}
\makeatother
\title{Base change for derived pushforward against a given map to K'}
\author{The Stacks Project Authors}
\date{stacks-project:spaces-perfect.tex:local-snapshot-20261012; GNU FDL 1.2 or later}
\begin{document}
\maketitle
\noindent\textit{Editorial note.} All delivered mathematical text is the author's own, copied from the pinned Stacks Project chapter file spaces-perfect.tex; no statement or proof step was rewritten, reordered or removed.\par\medskip
\noindent\textit{Editorial note.} The counted proof is the author's complete argument: the source proves the scheme case by the corresponding Derived Categories of Schemes lemma and carries out the induction step in full with an explicit commutative diagram, closing with the third-isomorphism lemma for distinguished triangles\par\medskip
\noindent\textit{Editorial note.} The label lemma-single-complex-base-change names the counted statement, and the proof environment immediately following it in the source is delivered verbatim as its proof.\par\medskip
\noindent\textit{Editorial note.} The counted statement is carried by the author's own lemma environment in the chapter source.\par\medskip
\noindent\textit{Allowed standard prerequisites (editorial).}\begin{itemize}
\item cartesian diagrams of algebraic spaces and the base change morphism Lg\textasciicircum{}*Rf\_*K to Rf'\_*L(g')\textasciicircum{}*K
\item the base change condition of the condition lemma (counted earlier in this lane) for a single complex K with a map L(g')\textasciicircum{}*K to K'
\item the induction principle for properties of algebraic spaces etale locally
\item relative Mayer-Vietoris triangles for distinguished squares
\end{itemize}\medskip
\ifdefined\hypertarget\hypertarget{sourcepointer-1}{}\fi
\noindent\textit{Source reference.} Original lemma 014A: sources/stacks-chow/tags.txt, source lines 1465--1465. Prerequisite pointer invoked at proof line(s) 5562; it names a standard fact of a sibling chapter that the counted argument uses, and every step of the counted proof built on it is retained verbatim in the delivered body.\par\medskip
\ifdefined\hypertarget\hypertarget{sourcepointer-2}{}\fi
\noindent\textit{Source reference.} Original reference 071Q: sources/stacks-spaces-perfect/pinned-original-source.tex, source lines 402--402. Prerequisite pointer invoked at proof line(s) 5525; it names a standard fact of the same chapter that the counted argument uses, and every step of the counted proof built on it is retained verbatim in the delivered body.\par\medskip
\ifdefined\hypertarget\hypertarget{sourcepointer-3}{}\fi
\noindent\textit{Source reference.} Original reference 08DR: sources/stacks-spaces-perfect/pinned-original-source.tex, source lines 1563--1563. Prerequisite pointer invoked at proof line(s) 5514; it names a standard fact of the same chapter that the counted argument uses, and every step of the counted proof built on it is retained verbatim in the delivered body.\par\medskip
\ifdefined\hypertarget\hypertarget{sourcepointer-4}{}\fi
\noindent\textit{Source reference.} Original reference 0DJ8: sources/stacks-spaces-perfect/pinned-original-source.tex, source lines 5365--5365. Prerequisite pointer invoked at proof line(s) 5503; it names a standard fact of the same chapter that the counted argument uses, and every step of the counted proof built on it is retained verbatim in the delivered body.\par\medskip
\ifdefined\hypertarget\hypertarget{sourcepointer-5}{}\fi
\noindent\textit{Source reference.} Original reference 08HZ: sources/stacks-spaces-perfect/pinned-original-source.tex, source lines 1979--1979. Prerequisite pointer invoked at proof line(s) 5548; it names a standard fact of the same chapter that the counted argument uses, and every step of the counted proof built on it is retained verbatim in the delivered body.\par\medskip
\ifdefined\hypertarget\hypertarget{sourcepointer-6}{}\fi
\noindent\textit{Source reference.} Original lemma 0DJ9: sources/stacks-chow/tags.txt, source lines 16969--16969. Prerequisite pointer invoked at proof line(s) 5536; it names a standard fact of a sibling chapter that the counted argument uses, and every step of the counted proof built on it is retained verbatim in the delivered body.\par\medskip
\ifdefined\hypertarget\hypertarget{sourcepointer-7}{}\fi
\noindent\textit{Source reference.} Original lemma 0DJ8: sources/stacks-chow/tags.txt, source lines 16968--16968. Prerequisite pointer invoked at proof line(s) 5528; it names a standard fact of a sibling chapter that the counted argument uses, and every step of the counted proof built on it is retained verbatim in the delivered body.\par\medskip
\ifdefined\hypertarget\hypertarget{sourcepointer-8}{}\fi
\noindent\textit{Source reference.} Original reference 08GH: sources/stacks-spaces-perfect/pinned-original-source.tex, source lines 896--896. Prerequisite pointer invoked at proof line(s) 5532; it names a standard fact of the same chapter that the counted argument uses, and every step of the counted proof built on it is retained verbatim in the delivered body.\par\medskip
\begin{lemma}
\label{lemma-single-complex-base-change}
Let $S$ be a scheme. Let
$$
\xymatrix{
X' \ar[r]_{g'} \ar[d]_{f'} &
X \ar[d]^f \\
Y' \ar[r]^g &
Y
}
$$
be a cartesian diagram of algebraic spaces over $S$.
Let $K \in D_\QCoh(\mathcal{O}_X)$ and let $L(g')^*K \to K'$
be a map in $D_\QCoh(\mathcal{O}_{X'})$. If
\begin{enumerate}
\item the equivalent conditions of
Lemma \ref{lemma-single-complex-base-change-condition} hold, and
\item $f$ is quasi-compact and quasi-separated,
\end{enumerate}
then the composition $Lg^*Rf_*K \to Rf'_*L(g')^*K \to Rf'_*K'$
is an isomorphism.
\end{lemma}

\begin{proof}
To check the map is an isomorphism we may work \'etale locally on $Y'$.
Hence we may assume $g : Y' \to Y$ is a morphism of affine schemes.
In this case, we will use the induction principle of
Lemma \ref{lemma-induction-principle}
to prove that for a quasi-compact and quasi-separated
algebraic space $U$ \'etale over $X$
the similarly constructed map
$Lg^*R(U \to Y)_*K|_U \to R(U' \to Y')_*K'|_{U'}$
is an isomorphism. Here $U' = X' \times_{g', X} U = Y' \times_{g, Y} U$.

\medskip\noindent
If $U$ is a scheme (for example affine), then the result holds.
Namely, then $Y, Y', U, U'$ are schemes, $K$ and $K'$ come from
objects of the derived category of the underlying schemes by
Lemma \ref{lemma-derived-quasi-coherent-small-etale-site}
and the condition of
Derived Categories of Schemes,
Lemma \ref{perfect-lemma-single-complex-base-change-condition}
holds for these complexes by
Lemma \ref{lemma-single-complex-base-change-condition}.
Thus (by the compatibilities explained in
Remark \ref{remark-match-total-direct-images})
we can apply the result in the case of schemes
which is
Derived Categories of Schemes, Lemma
\ref{perfect-lemma-single-complex-base-change}.

\medskip\noindent
The induction step. Let $(U \subset W, V \to W)$ be an elementary
distinguished square with $W$ a quasi-compact and quasi-separated
algebraic space \'etale over $X$, with $U$ quasi-compact, $V$ affine
and the result holds for $U$, $V$, and $U \times_W V$.
To easy notation we replace $W$ by $X$ (this is permissible at this point).
Denote $a : U \to Y$, $b : V \to Y$, and $c : U \times_X V \to Y$
the obvious morphisms. Let $a' : U' \to Y'$, $b' : V' \to Y'$
and $c' : U' \times_{X'} V' \to Y'$ be the base changes of $a$, $b$, and $c$.
Using the distinguished triangles from relative Mayer-Vietoris
(Lemma \ref{lemma-unbounded-relative-mayer-vietoris})
we obtain a commutative diagram
$$
\xymatrix{
Lg^*Rf_*K \ar[r] \ar[d] & Rf'_*K' \ar[d] \\
Lg^*Ra_*K|_U \oplus Lg^*Rb_*K|_V \ar[r] \ar[d] &
Ra'_* K'|_{U'} \oplus Rb'_* K'|_{V'} \ar[d] \\
Lg^*Rc_*K|_{U \times_X V} \ar[r] \ar[d] &
Rc'_*K'|_{U' \times_{X'} V'} \ar[d] \\
Lg^*Rf_* K[1] \ar[r] &
Rf'_* K'[1]
}
$$
Since the 2nd and 3rd horizontal arrows are isomorphisms so is the first
(Derived Categories, Lemma \ref{derived-lemma-third-isomorphism-triangle})
and the proof of the lemma is finished.
\end{proof}

\end{document}
